\begin{papercard}{hutton2018}
\cf{Stages} \stagetag{D} \stagetag{W} \stagetag{C} (also \stagetag{A}, \stagetag{I}). A privacy \emph{heuristic evaluation} instrument for self-tracking apps, applied to 64 Android apps with a partial double-rating check.

\cf{Aim} To understand how mass-market self-tracking apps perform on privacy, by introducing heuristics for privacy characteristics and applying them to 64 popular services \pc{hutton2018}{1}{introducing a set of heuristics for evaluating privacy characteristics}. Secondary question: do app categories, or mHealth vs.\ non-health self-tracking, differ in privacy performance?

\cf{Scope of what is evaluated} Only what users can know or control \emph{at the point of use}: the interface, terms of service and privacy policy. The ``invisible'' facets, such as sharing or selling without consent, are deliberately excluded because they cannot be assessed when deciding to use an app \pc{hutton2018}{2}{do not focus on the invisible facets of privacy}. This is the same communication-only boundary as in the thesis.

\cf{Instrument development} Four user questions frame the heuristics: Am I informed before use? Do I have control once using it? Do I have access to my data? Can I control disclosure to third parties? \pc{hutton2018}{3}{Am I informed about what happens to my data before I use}. Sources: three FTC Fair Information Practice categories (notice/awareness, choice/consent, access/participation) as the guiding frame, self-tracking privacy concerns from the literature, GDPR provisions (portability, unambiguous consent, erasure), the STRAP framework, and Inostroza et al.'s touchscreen usability heuristics for a fourth dimension, \emph{social disclosure usability} \pc{hutton2018}{3}{categories from the FTC's Fair Information Practices}. Three expert reviewers (HCI/software engineering, personal self-tracking experience) pilot-tested the set for consistent application and exhaustiveness. Four heuristics were rearranged and one was added, giving 26 heuristics in total \pc{hutton2018}{3}{An initial set of heuristics was pilot-tested by 3 reviewers}. Items are dichotomous or ordinal, with higher scores meaning more privacy-preserving \pc{hutton2018}{3}{The heuristics include a}.

\cf{The 26 heuristics} Notice (H1--H7): entity, uses, recipients and nature of data identified before remote sharing; safeguards explained; sufficiently explicit notice; control over secondary use. Choice/consent (H8--H15): consent before sharing; explicit opt-in; choice of data types; dynamic consent for new data, new purposes and new recipients; revocability; control over storage location. Access (H16--H19): full raw export; standard formats; in-service export; API access. Social disclosure usability (H20--H26): per-disclosure and granular controls, confirmation, memory load, minimalism, consistency, contextual help \pc{hutton2018}{4}{The 26 heuristics (H) present in our evaluation technique}. The scoring levels are only in the separate scoring sheet \pp{hutton2017heuristics}. Example: H6 runs from 0 (no data usage explained) through 2 (shown in situ and must be scrolled through) to 4 (must be read and comprehension is assessed) \pc{hutton2017heuristics}{1}{For those of above satisfied, notice is sufficiently explicit}. ``No remote sharing'' earns the maximum on the notice items \pc{hutton2017heuristics}{1}{Entity is identified 2: No remote sharing}. H14 grades revocation from 0 (none) through 1 (contact support) to 3 (all actors revocable in-app) \pc{hutton2017heuristics}{3}{Consent to store and process data can be revoked at any time}.

\cf{Sampling frame and search} UK Google Play Store, Android 6.0 test phone. Android was chosen for market share and because Google's less restrictive vetting yields a broad range of apps \pc{hutton2018}{5}{We focused on this platform because the Android mobile operating}. Search categories were derived from a 2014 survey of 105 London Quantified Self Meetup members, because no taxonomy of self-tracking apps existed \pc{hutton2018}{5}{survey of 105 members of the London Quantified Self Meetup}. The 20 most frequently reported data types became keywords (e.g.\ weight, sleep, mood), run through a scraping script that saved store metadata to CSV \pc{hutton2018}{5}{the 20 most-frequently reported types of data to produce a series}. Keywords were grouped into 12 categories, mHealth vs.\ other self-tracking (Table 2) \pc{hutton2018}{5}{The categories of apps used in our evaluation and the constituent search keywords}.

\cf{Eligibility and selection} 345 hits, of which instructional exercise apps were excluded because they capture no data, leaving 292 \pc{hutton2018}{5}{Our search produced an initial list of 345 apps}. Two major vendors missing from the search (Xiaomi Mi Fit, Jawbone UP) were added manually from market data, and code and dataset are public \pc{hutton2018}{5}{to identify major commercial}. The seven most popular apps per category (fewer where unavailable) gave 64 apps \pc{hutton2018}{5}{selecting the 7 most popular apps}. The app list with package name, version tested, matching keywords and download band is in the second appendix \pp{hutton2017heuristics}.

\cf{Procedure} Each app was assigned to one of four reviewers, who checked the user interface, terms of service and privacy policy against all 26 heuristics \pc{hutton2018}{5}{Reviewers were instructed to review the user}. Inapplicable heuristics were scored ``not applicable'' and excluded from that app's denominator (e.g.\ H20--H26 for apps without social sharing). Applicability is at the evaluator's discretion \pc{hutton2018}{5}{if a heuristic did not apply}.

\cf{Reliability} Twelve random apps were rated by a second reviewer, allocated so that every reviewer checked every peer at least once. Cohen's $\kappa=.45$ (moderate). Disagreements came mainly from unclear policy language and were resolved by consensus \pc{hutton2018}{6}{To assess the interrater reliability in applying the heuristics, 12}. The authors attribute the moderate agreement to varying interpretation of legal language and to double-rating only a subset, and suggest more raters per app \pc{hutton2018}{12}{Cohen kappa for our interrater reliability}.

\cf{Analysis} Overall privacy score = sum over applicable heuristics, also expressed as a percentage of the attainable maximum \pc{hutton2018}{6}{The outcome of the analysis was an overall privacy score}. Cumulative link mixed models (ordinal regression with random effects for repeated ratings) compared categories, dimensions and heuristics \pc{hutton2018}{6}{we fit an appropriate cumulative link mixed model}.

\cf{Key findings} Mean satisfaction 46.2\,\% (SD 24.3). App category was not a significant predictor, and popularity proxies (downloads, stars, number of ratings) did not predict privacy \pc{hutton2018}{6}{attempts to explore this question using the number of downloads}. Access was the weakest dimension, with paywalled export scored as unsatisfactory \pc{hutton2018}{6}{We considered a paywall between users and their own data}. mHealth apps scored significantly worse than non-health self-tracking apps \pc{hutton2018}{7}{that nonhealth-related self-tracking apps performed significantly}. Of 51 apps collecting personal data, only 29 had terms or a policy, and none required reading before registration \pc{hutton2018}{10}{only 29 included terms of service or a privacy}. Social sign-on often bypassed the terms and marketing consent \pc{hutton2018}{11}{such a registration route bypassed the usual exposure to an app}.

\cf{Limitations stated} Survey-derived keywords may be biased, so some apps were possibly missed. The most popular app, a period calendar, entered only via the ``mood'' keyword \pc{hutton2018}{12}{the most popular app that we examined was Simple Design}. UK store only, so results are region-specific \pc{hutton2018}{12}{Our study used the UK Google Play Store}. The heuristics were applied context-free by experts. Application is time-consuming because it needs close reading of the terms and exploration of many app routes \pc{hutton2018}{12}{the heuristics can be time consuming}.

\cf{Use for this thesis} The strongest template for an \emph{observable checklist with ordinal levels} and an explicit ``not applicable'' rule with adjusted denominators. The notice/consent/access items map onto walkthrough tasks (consent, deletion, export). Its reliability design (second rater on a random subset, kappa, then consensus) is a minimal precedent; the thesis should report per-code agreement on frozen labels. Its finding that popularity does not predict privacy is useful when justifying a popularity-based sample: popularity measures exposure, not quality.
\end{papercard}

\begin{papercard}{hutton2017heuristics}
\cf{What it is} The operational scoring sheet behind the 26 heuristics of Hutton et al.\ (2018), version 1 (23 Oct 2017). Each heuristic has a numbered range and anchored levels \pc{hutton2017heuristics}{1}{Quantified self privacy heuristics}.
\cf{Structure} Notice/awareness H1--H7 (H1--H5 scored 0--2, with 2 = no remote sharing; H6 0--4; H7 0--2); choice/consent H8--H15 (H12--H14 up to 3) \pc{hutton2017heuristics}{2}{Data processing consent is dynamic}; access/participation H16--H19 (H16 0--3: none, aggregate, subset, all raw) \pc{hutton2017heuristics}{3}{All raw collected data can be extracted from the service}; social disclosure usability H20--H26 \pc{hutton2017heuristics}{4}{Social disclosure usability}.
\cf{Use for this thesis} Copy the \emph{format}: anchored ordinal levels that describe observable states, e.g.\ ``deletion only by contacting support'' vs.\ ``in-app''. This makes two-coder agreement measurable and reduces interpretive drift. Example: H18 distinguishes extraction by contacting the provider (1) from extraction at any time within the service (2) \pc{hutton2017heuristics}{3}{Data can be extracted by contacting the service provider}.
\end{papercard}
